\heading{PHỤ LỤC}

Các phụ lục trích ý chính của mã nguồn và cấu hình; phần lược bỏ được đánh dấu bằng dấu ba chấm. Toàn bộ mã nguồn nằm trong kho mã ở Phụ lục A.

\subsection*{Phụ lục A. Mã nguồn và dữ liệu}
\addcontentsline{toc}{subsection}{Phụ lục A. Mã nguồn và dữ liệu}
\textit{A.1. Kho mã}

Toàn bộ cấu hình, mã mock, kiểm thử và bằng chứng của đề tài nằm trong kho mã công khai \url{https://github.com/QuanVongHuu/MCP-gateway-rbac}. Nội dung báo cáo được đối chiếu với nhánh \texttt{docs/final\_thesis} tại commit \texttt{7e4aa42} (ngày 05/11/2026).
% Commit 7e4aa42 chưa được push lên GitHub lúc viết; cập nhật hash sau lần push cuối cùng.

\textit{A.2. Cấu trúc thư mục}

{\singlespacing\renewcommand{\arraystretch}{1.25}\noindent
\begin{tabularx}{\textwidth}{|p{3.4cm}|X|}\hline
\textbf{Thư mục} & \textbf{Nội dung} \\ \hline
\texttt{deploy/} & Manifest Kubernetes theo từng lớp: \texttt{agentgateway/} (Gateway, HTTPRoute, backend năm target), \texttt{keycloak/}, \texttt{auth/} (xác thực JWT), \texttt{identity/} (token exchange, Kubernetes RBAC), \texttt{network/} (NetworkPolicy), \texttt{cluster/} và \texttt{real-servers/} (radar, OpenMetadata, Spark History thật) \\ \hline
\texttt{rbac/} & \texttt{tool-access.yaml}, nguồn sự thật duy nhất của quy tắc vai trò -- công cụ; ghi chép thí nghiệm về các chính sách chồng nhau \\ \hline
\texttt{servers/} & Năm mock MCP server viết bằng FastMCP, mỗi thư mục mang đúng tên target; module chung \texttt{mockkit.py}, \texttt{Dockerfile} và \texttt{k8s/mocks.yaml} \\ \hline
\texttt{tests/} & Kiểm thử đầu cuối \texttt{gateway\_smoke.py}, \texttt{agent\_rbac\_smoke.py}; kiểm thử tĩnh \texttt{static\_contract\_test.mjs}; ma trận kỳ vọng trong \texttt{expected/}; kịch bản \texttt{curl/} và \texttt{inspector/} \\ \hline
\texttt{scripts/} & \texttt{up.sh} áp các lớp cấu hình theo thứ tự phụ thuộc; các script lấy token, sinh secret cho agent và chạy kiểm thử \\ \hline
\texttt{ops/} & \texttt{up.ps1} và \texttt{down.ps1} cho Windows PowerShell \\ \hline
\texttt{observability/} & Trường bổ sung cho access log (\texttt{access-logs/}); Alloy, Loki, Grafana và dashboard allow/deny (\texttt{grafana/}) \\ \hline
\texttt{evidence/} & Đầu ra kiểm thử, trích đoạn log và số đo làm bằng chứng cho Chương 5 \\ \hline
\texttt{docs/}, \texttt{Kt/} & Ghi chú thiết kế, kế hoạch triển khai, ma trận quyền của AI Agent; sơ đồ kiến trúc draw.io và mô tả luồng \\ \hline
\texttt{report/} & Mã nguồn LaTeX của báo cáo này \\ \hline
Thư mục gốc & \texttt{bootstrap.sh}, \texttt{run.sh}, \texttt{QUICKSTART.md}, \texttt{AGENTS.md} \\ \hline
\end{tabularx}}

\textit{A.3. Chạy nhanh}

Môi trường cần kind, kubectl, helm, Docker có Compose (khoảng 8 GB bộ nhớ), Node 22.19 trở lên, curl, jq và Python có gói \texttt{mcp==1.30.0}. Hướng dẫn đầy đủ nằm trong \texttt{QUICKSTART.md}.
\begin{lstlisting}
# once (Git Bash or WSL)
./bootstrap.sh     # kind cluster, Gateway API CRDs, agentgateway v1.5.0, mock image
# every start
./run.sh           # K8s layers, real backends, Grafana, port-forwards 8080/8002
.\ops\up.ps1       # Windows: scripts/up.sh + Grafana, run in WSL (Ubuntu)
# checks
./run.sh test                         # run.sh, then scripts/run-tests.sh
bash scripts/run-agent-rbac-tests.sh  # tests/agent_rbac_smoke.py
node tests/static_contract_test.mjs   # manifests vs. matrix, no cluster
# cleanup
.\ops\down.ps1     # delete the kind cluster
\end{lstlisting}
\texttt{ops/up.ps1} không khởi động ba backend thật. Khi agent mở phiên, gateway kết nối tới mọi target, nên chỉ cần một target ngừng là mọi agent đều lỗi; vì vậy script này chỉ dùng được khi các backend đó đã chạy sẵn. \texttt{ops/down.ps1} chỉ xóa cụm kind, còn hai nhóm container Docker Compose vẫn chạy.

\textit{A.4. Ghi chú về dữ liệu}

Mọi dữ liệu trong kho là dữ liệu thử. Năm mock trả dữ liệu dựng sẵn cho đủ 21 công cụ; các công cụ mang tính ghi như \texttt{airflow\_trigger\_dag} hay \texttt{radar\_manage\_workload} chỉ trả kết quả mô phỏng, không gây tác dụng phụ. Cấu hình hiện tại trỏ ba target radar, OpenMetadata và Spark History sang phần mềm thật chạy cục bộ, nhưng các hệ thống này cũng chỉ chứa dữ liệu thử: radar đọc chính cụm kind của hệ thống, OpenMetadata có người dùng thử và một bảng mẫu do script cài đặt tạo ra, Spark History đọc ba event log mẫu của dự án gốc. Kho không chứa dữ liệu nghiệp vụ của đơn vị kiến tập.

Thư mục \texttt{evidence/} chỉ lưu đầu ra đã lược bỏ bí mật, không có client secret, access token hay claim nhạy cảm. Secret của từng agent do \texttt{scripts/\allowbreak bootstrap-agent-secrets.sh} sinh thẳng vào Kubernetes Secret mà không in ra màn hình; tệp cấu hình realm chỉ chứa biến giữ chỗ như \texttt{\$\{CLUSTER\_OBSERVER\_AGENT\_SECRET\}}. Tài khoản \texttt{admin}/\texttt{admin} của Keycloak và Grafana cùng mật khẩu \texttt{password} của người dùng thử là giá trị mặc định, chỉ dành cho môi trường thử cục bộ.

\subsection*{Phụ lục B. Cấu hình gateway}
\addcontentsline{toc}{subsection}{Phụ lục B. Cấu hình gateway}
\textit{B.1. Backend hợp nhất năm target} (\texttt{deploy/\allowbreak agentgateway/\allowbreak AgentgatewayBackend.yaml})
\begin{lstlisting}
kind: AgentgatewayBackend
spec:
  mcp:
    prefixMode: Always          # tools appear as {target}_{tool}
    targets:
    - name: dbt
      static: {host: dbt.mcp.svc.cluster.local, port: 80, path: /mcp,
               protocol: StreamableHTTP}
    ...
\end{lstlisting}
\textit{B.2. Xác thực JWT và chuyển tiếp định danh} (\texttt{deploy/auth/jwt.yaml})
\begin{lstlisting}
traffic:
  jwtAuthentication:
    mode: Strict
    preserveToken: true
    providers:
    - issuer: http://localhost:8002/realms/mcp
      audiences: [mcp-gateway]
      jwks: {remote: {jwksPath: /realms/mcp/protocol/openid-connect/certs,
                      cacheDuration: "5m", ...}}
  transformation:
    request:
      set:
      - {name: x-forwarded-user,   value: "jwt.preferred_username"}
      - {name: x-agent-id,         value: "has(jwt.azp) ? jwt.azp : ''"}
      - {name: x-forwarded-groups, value: "has(jwt.roles) ? jwt.roles.join(',') : ''"}
      - {name: x-caller-token,     value: 'request.headers["authorization"]'}
      remove: [authorization]
---
# second policy, on the backend: drop the token before it leaves the gateway
backend: {transformation: {request: {remove: [x-caller-token]}}}
\end{lstlisting}
\textit{B.3. Phân quyền theo công cụ} (\texttt{rbac/tool-access.yaml})
\begin{lstlisting}
backend:
  mcp:
    authorization:
      action: Allow
      policy:
        matchExpressions:      # expressions are OR-ed
        - 'has(jwt.roles) && "analyst" in jwt.roles
           && mcp.tool.target == "openmetadata"
           && mcp.tool.name in ["search_metadata", "get_entity_details", ...]'
        ...
        - 'has(jwt.roles) && "kube_operator" in jwt.roles
           && mcp.tool.target == "radar" && mcp.tool.name == "manage_workload"'
        - 'has(jwt.roles) && "platform_admin" in jwt.roles'
\end{lstlisting}
\textit{B.4. Cô lập mạng} (\texttt{deploy/network/backends.yaml})
\begin{lstlisting}
kind: NetworkPolicy
metadata: {name: backends-from-gateway-only, namespace: mcp}
spec:
  podSelector: {}               # every pod in namespace mcp
  policyTypes: [Ingress]
  ingress:
  - from:
    - namespaceSelector:
        matchLabels: {kubernetes.io/metadata.name: agentgateway-system}
\end{lstlisting}

\subsection*{Phụ lục C. Mock MCP server}
\addcontentsline{toc}{subsection}{Phụ lục C. Mock MCP server}
\textit{C.1. Bộ khung chung} (\texttt{servers/mockkit.py})
\begin{lstlisting}[language=Python]
def new_server(name):
    # stateless + JSON: any gateway replica can serve any request
    return FastMCP(name, host="0.0.0.0", port=8000,
                   stateless_http=True, json_response=True)

def seen(ctx, target, tool, args):
    """Log who called which tool. Never logs the token itself."""
    headers = ctx.request_context.request.headers
    record = {"target": target, "tool": tool, "args": args,
              "x_agent_id": headers.get("x-agent-id"),
              "x_forwarded_user": headers.get("x-forwarded-user"),
              "x_forwarded_groups": headers.get("x-forwarded-groups"),
              "authorization_present": "authorization" in headers}
    if token := bearer(headers):          # non-sensitive claims only
        claims = decode(token)
        record["token"] = {k: claims.get(k) for k in
                           ("iss", "aud", "azp", "sub", "preferred_username", "roles")}
    print(json.dumps(record))
    return record
\end{lstlisting}
\textit{C.2. Một công cụ mẫu} (\texttt{servers/dbt/server.py})
\begin{lstlisting}[language=Python]
@mcp.tool()
def test(ctx, selector=""):
    """EXECUTE: run dbt data tests. Mock: nothing is run."""
    who = seen(ctx, "dbt", "test", {"selector": selector})
    return {"result": "mock: dbt test -> PASS=17 WARN=1 ERROR=1",
            "seen_by_mock": who}
\end{lstlisting}
\textit{C.3. Chức năng của 21 công cụ}
{\singlespacing\renewcommand{\arraystretch}{1.2}\small
\begin{longtable}{|p{2.4cm}|p{3.8cm}|c|p{6.4cm}|}\hline
\textbf{Target} & \textbf{Công cụ} & \textbf{Loại} & \textbf{Chức năng} \\ \hline\endhead
radar & list\_namespaces & Đọc & Liệt kê namespace của cụm \\ \hline
radar & list\_resources & Đọc & Tóm tắt tài nguyên theo loại \\ \hline
radar & get\_pod\_logs & Đọc & Lấy log gần nhất của pod \\ \hline
radar & manage\_workload & Ghi & Khởi động lại, scale, rollback workload \\ \hline
spark-history & list\_applications & Đọc & Liệt kê ứng dụng Spark theo trạng thái \\ \hline
spark-history & get\_stage & Đọc & Chi tiết một stage và chỉ số task \\ \hline
spark-history & get\_job\_bottlenecks & Đọc & Tìm stage chậm, task lệch, lỗi executor \\ \hline
openmetadata & search\_metadata & Đọc & Tìm tài sản dữ liệu theo từ khóa, owner, tag \\ \hline
openmetadata & get\_entity\_details & Đọc & Chi tiết một thực thể theo tên đầy đủ \\ \hline
openmetadata & get\_entity\_lineage & Đọc & Nguồn phía trên và đích phía dưới \\ \hline
openmetadata & patch\_entity & Ghi & Sửa mô tả, tag, owner bằng JSON Patch \\ \hline
dbt & list\_metrics & Đọc & Liệt kê metric của semantic layer \\ \hline
dbt & query\_metrics & Đọc & Truy vấn metric theo chiều \\ \hline
dbt & get\_model\_health & Đọc & Trạng thái chạy, test, độ tươi của model \\ \hline
dbt & get\_lineage & Đọc & Cha và con của một node \\ \hline
dbt & test & Thực thi & Chạy data test theo selector \\ \hline
airflow & list\_dags & Đọc & Liệt kê DAG và lần chạy gần nhất \\ \hline
airflow & get\_dag\_run & Đọc & Trạng thái một lần chạy DAG \\ \hline
airflow & get\_task\_logs & Đọc & Log của một lần thử task \\ \hline
airflow & trigger\_dag & Ghi & Kích hoạt một lần chạy mới \\ \hline
airflow & pause\_dag & Ghi & Tạm dừng DAG \\ \hline
\end{longtable}}

\subsection*{Phụ lục D. Danh tính AI Agent}
\addcontentsline{toc}{subsection}{Phụ lục D. Danh tính AI Agent}
\textit{D.1. Định nghĩa một client agent} (\texttt{deploy/keycloak/realm-config.yaml})
\begin{lstlisting}
{"clientId": "cluster-observer-agent",
 "publicClient": false, "serviceAccountsEnabled": true,
 "standardFlowEnabled": false, "directAccessGrantsEnabled": false,
 "secret": "${CLUSTER_OBSERVER_AGENT_SECRET}",
 "protocolMappers": [ audience mcp-gateway, client roles -> claim "roles",
                      hardcoded claim kube_group = "kube_viewer" ]}
...
"users": [{"serviceAccountClientId": "cluster-observer-agent",
           "clientRoles": {"mcp-gateway": ["kube_viewer"]}}, ...]
\end{lstlisting}
\textit{D.2. Nạp agent vào realm} (\texttt{deploy/keycloak/setup-agents.sh})
\begin{lstlisting}[language=bash]
for cid in <clients with a secret>; do   # replace placeholders with real secrets
  sec=$(kubectl get secret agent-client-secrets -o jsonpath={.data.$cid} | base64 -d)
  json=$(jq '(.clients[] | select(.clientId == $c)).secret = $s' ...)
done
# partial import: agent clients and roles only, mcp-gateway untouched
kcadm.sh create partialImport -r mcp  (ifResourceExists: OVERWRITE)
kcadm.sh add-roles --uusername service-account-$cid --rolename $role
\end{lstlisting}

\subsection*{Phụ lục E. Tự động hóa triển khai}
\addcontentsline{toc}{subsection}{Phụ lục E. Tự động hóa triển khai}
\textit{E.1. Lớp bọc PowerShell} (\texttt{ops/up.ps1})
\begin{lstlisting}
$repoDir = (& wsl.exe -d $Distro -- wslpath -u $PSScriptRoot).Trim()
$command = "cd '$repoDir' && bash scripts/up.sh && bash observability/grafana/up.sh"
$stopwatch = [System.Diagnostics.Stopwatch]::StartNew()
& wsl.exe -d $Distro -- bash -lc $command
if ($LASTEXITCODE -ne 0) { throw "Startup failed after $($stopwatch.Elapsed)" }
\end{lstlisting}
\textit{E.2. Áp các lớp cấu hình} (\texttt{scripts/up.sh}, rút gọn)
\begin{lstlisting}[language=bash]
set -euo pipefail
kubectl apply -f deploy/agentgateway/Gateway.yaml -f servers/k8s/mocks.yaml ...
kubectl apply -f deploy/keycloak/keycloak.yaml
# setup-realm.sh rebuilds the realm, so only run it when the realm is missing
kcadm.sh get realms/mcp || bash deploy/keycloak/setup-realm.sh
bash scripts/bootstrap-agent-secrets.sh && bash deploy/keycloak/setup-agents.sh
kubectl apply -f deploy/auth/jwt.yaml -f deploy/identity/openmetadata.yaml
kubectl apply -f rbac/tool-access.yaml
kubectl apply -f deploy/network/backends.yaml
kubectl apply -f observability/access-logs/policy.yaml
\end{lstlisting}
\textit{E.3. Script dọn dẹp} (\texttt{ops/down.ps1})
\begin{lstlisting}
& kind delete cluster --name $ClusterName
if ($LASTEXITCODE -ne 0) { Write-Error "Failed to delete ..."; exit $LASTEXITCODE }
\end{lstlisting}

\subsection*{Phụ lục F. Giám sát allow/deny}
\addcontentsline{toc}{subsection}{Phụ lục F. Giám sát allow/deny}
\textit{F.1. Bổ sung trường vào access log} (\texttt{observability/\allowbreak access-logs/\allowbreak policy.yaml})
\begin{lstlisting}
frontend:
  accessLog:
    attributes:
      add:
      - {name: user,       expression: 'has(jwt.preferred_username) ? ... : "anonymous"'}
      - {name: agent_id,   expression: 'has(jwt.azp) ? jwt.azp : ""'}
      - {name: roles,      expression: 'has(jwt.roles) ? jwt.roles : []'}
      - {name: mcp_method, expression: 'has(mcp.methodName) ? mcp.methodName : ""'}
      - {name: mcp_target, expression: 'has(mcp.tool.target) ? mcp.tool.target : ""'}
      - {name: mcp_tool,   expression: 'has(mcp.tool.name) ? mcp.tool.name : ""'}
\end{lstlisting}
\textit{F.2. Truy vấn LogQL đếm deny theo vai trò} (\texttt{observability/\allowbreak grafana/\allowbreak dashboard.yaml})
\begin{lstlisting}
sum by (roles) (count_over_time(
  {namespace="agentgateway-system"}
  |~ `mcp.method.name=tools/call` |~ `http.status=400` |~ `Unknown tool`
  | regexp `roles=(?P<roles>\[[^]]*\])` [$__interval]))
\end{lstlisting}
